% !TEX root = ../../cookbook.tex
\newpage
\recipe{Mac \& Cheese}{Daniel's favorite. Serves 4 (o 1.5 Daniels)}{}{}

\subsection*{Ingredients}
\begin{ingredients}
	\item 460 g elbow macaroni
	\item 60 g butter
	\item 36 g AP flour
	\item 480 ml milk
	\item Grated cheese: 
	\begin{ingredients}
		\item 70 g blue cheese
		\item 70g parmesan
		\item 80 g sharp cheddar
		\item 60 g gouda
	\end{ingredients}
	\item 4 cloves of garlic, \sfrac{1}{2} crushed, and \sfrac{1}{2} finely chopped
	\item Salt and pepper
	\item 1 Habanero, or other pepper of choice, finely diced (optional) 
	\item Olive oil (optional)
\end{ingredients}

\medskip

% --- Preparation ---
Start cooking the \textbf{macaroni} to \textit{al dente}.
In parallel, if you like a pleasurably spicy-hot mac, pan fry the \textbf{Habanero} in a separate pan with a tablespoon of olive oil, until slightly golden. 

In the meantime, melt the \textbf{butter} in a large stainless steel pan, add the \textbf{flour} and vigorously whisk until smooth\footnote{The result is what the French call ``\textit{roux}''}. 
After the buttery flour mix has been bubbling for a couple of minutes, add the crushed garlic\footnote{For a stronger garlic taste, add less garlic at this stage, leaving more to be added raw later.} and mix for another minute.
While whisking, slowly add the \textbf{milk} to the mixture.
Progressively add all the \textbf{grated cheese} and, finally add the diced \textbf{garlic}, just before taking the sauce off the heat.

Season to taste with \textbf{salt and pepper}, then serve on the warm macaroni.

\vspace*{0.75 cm}

% --- Description ---
\begin{recipeblock}
	
\end{recipeblock}